\chapter{第二讲补编}

\section*{5.0.1 波映射}

从 Minkowski 时空到 Riemann 流形 $(M,g)$ 的波映射, 是泛函
\begin{equation}
\mathcal{L}[u]=\frac12\int_{\mathbb{R}^{1+n}}\langle du,du\rangle\,dt\,dx
\end{equation}
关于具紧支集变分的临界点, 其中在 $M$ 的局部坐标下 $\langle du,du\rangle=\sum_{\mu=0}^n\sum_{a,b}g_{ab}\partial_\mu u^a\partial^\mu u^b$. 该变分问题的 Euler--Lagrange 方程恰为 (WM) 的形式(在 $M$ 的局部坐标下), 其中 $\Gamma_{JK}^I$ 为 $M$ 在局部坐标系下的 Christoffel 符号, $N=\dim M$(参见例如 Shatah--Struwe \cite{klai58}).
\section*{5.0.2 Maxwell--Klein--Gordon 方程}

在下面的讨论中采用求和约定: 希腊指标从 $0$ 到 $n$ 求和, 罗马指标从 $1$ 到 $n$ 求和. 注意指标是相对 Minkowski 度量 $m_{\mu\nu}=\mathrm{diag}(-1,1,\ldots,1)$ 升降的. 例如 $\Box=\partial^\mu\partial_\mu$, $\Delta=\partial^j\partial_j$. 记 $\mathrm{i}$ 为虚数单位.
方程的未知量是定义于 Minkowski 时空上的一个一次形式 $A_\mu dx^\mu$(规范势)与一个标量 $\phi$:
\begin{equation}
\begin{aligned}
A_\mu&:\mathbb{R}^{1+n}\to\mathbb{R},\\
\phi&:\mathbb{R}^{1+n}\to\mathbf{C}.\end{aligned}
\end{equation}
电磁场为二次形式 $F_{\mu\nu}=\partial_\mu A_\nu-\partial_\nu A_\mu$. 相对规范势的协变导数为
\begin{equation}
D_\mu\phi=\partial_\mu\phi+\mathrm{i}A_\mu\phi.\end{equation}
我们寻找泛函
\begin{equation}
\mathcal{L}[A_\mu,\phi]=\int_{\mathbb{R}^{1+n}}\Bigl(-\frac14 F_{\mu\nu}F^{\mu\nu}-\frac12 D_\mu\phi\,\overline{D^\mu\phi}\Bigr)\,dt\,dx
\end{equation}
的临界点. 相应的 Euler--Lagrange 方程为
\begin{gather}
\partial^\mu F_{\mu\nu}=-\Im\bigl(\phi\,\overline{D_\nu\phi}\bigr),\\
D^\mu D_\mu\phi=0,\end{gather}
其中 $\Im z$ 表示 $z$ 的虚部.
设 $\chi$ 为 $\mathbb{R}^{1+n}$ 上实值函数, 考虑变换 $(A_\mu,\phi)\to(\widetilde A_\mu,\widetilde\phi)$, 由
\begin{gather}
\widetilde A_\mu=A_\mu-\partial_\mu\chi,\\
\widetilde\phi=e^{\mathrm{i}\chi}\phi
\end{gather}
给出. 显然, 电磁场在规范变换 $A_\mu\to\widetilde A_\mu$ 下保持不变; 且易见若 $(A_\mu,\phi)$ 满足 (MKG), 则 $(\widetilde A_\mu,\widetilde\phi)$ 亦满足(注意 $D_\mu$ 依赖于 $A_\mu$). 这就在满足 (MKG) 的数对 $(A_\mu,\phi)$ 的集合上给出了一个等价关系; 所谓 (MKG) 的一个解, 我们理解为此类数对的等价类.
于是我们拥有规范自由度; 即我们可自由选取给定解(等价类)的任一代表, 并可以规定规范势所应满足的条件. 传统的规范条件有:
\begin{itemize}
\item Lorentz 规范: $\partial^\mu A_\mu=0$;
\item Coulomb 规范: $\partial^j A_j=0$;
\item 时间规范: $A_0=0$.\end{itemize}

\paragraph*{(MKG) 的 Lorentz 规范.} 将 Lorentz 条件与 (MKG) 结合, 得到系统
\begin{gather}
\Box A_\mu=-\Im(\phi\,\overline{\partial_\mu\phi})+|\phi|^2A_\mu,\\
\Box\phi=-2\mathrm{i}A^\mu\partial_\mu\phi+A^\mu A_\mu\phi,\\
\partial^\mu A_\mu=0.\end{gather}
现注意, 若 $(A_\mu,\phi)$ 满足 (110a) 与 (110b), 其初数据
\begin{gather}
A_\mu|_{t=0}=a_\mu,\qquad \partial_t A_\mu|_{t=0}=b_\mu,\\
\phi|_{t=0}=\phi_0,\qquad \partial_t\phi|_{t=0}=\phi_1
\end{gather}
满足约束
\begin{equation}
b_0=\partial^j a_j,\qquad \Delta a_0-|\phi_0|^2 a_0=\partial^j b_j-\Im(\phi_0\overline{\phi}_1),\end{equation}
则 (110c) 自动满足. 因为由 (110a) 与 (110b), 量 $u=\partial^\mu A_\mu$ 满足 $\Box u=|\phi|^2u$, 且由 (111a) 与 (112) 有 $u|_{t=0}=\partial_t u|_{t=0}=0$. 由解的唯一性知 $u=0$.
因此, (110c) 等价于初数据上的约束 (112), 于是我们只需处理 (110a) 与 (110b). 从而 Lorentz 规范下的 (MKG) 在形式上即为 $\Box u=u\partial u+u^3$. 遗憾的是, 此类一般方程并不具备良好的局部正则性, 故 Lorentz 规范对我们的目的并不十分有用.
\paragraph*{(MKG) 的 Coulomb 规范.} 将 Coulomb 条件与 (MKG) 结合, 得到
\begin{gather}
\Delta A_0=-\Im(\phi\,\overline{\partial_t\phi})+|\phi|^2A_0,\\
\Box A_j=-\Im(\phi\,\overline{\partial_j\phi})+|\phi|^2A_j-\partial_j\partial_t A_0,\\
\Box\phi=-2\mathrm{i}A^j\partial_j\phi+2\mathrm{i}A_0\partial_t\phi+\mathrm{i}(\partial_t A_0)\phi+A^\mu A_\mu\phi,\\
\partial^j A_j=0.\end{gather}
此处我们将规范势分解为时间分量 $A_0$ 与空间分量 $A=A_j dx^j$. 在 $t=0$ 给定初数据:
\begin{gather}
A_j|_{t=0}=a_j,\qquad \partial_t A_j|_{t=0}=b_j,\\
\phi|_{t=0}=\phi_0,\qquad \partial_t\phi|_{t=0}=\phi_1.\end{gather}
对 $A_0$ 不加初值条件; 若记 $a_0=A_0|_{t=0}$, 则由 (113a) 有 $\Delta a_0-|\phi_0|^2a_0=-\Im(\phi_0\overline{\phi}_1)$.
若数据散度自由, 即
\begin{equation}
\partial^j a_j=\partial^j b_j=0,\end{equation}
则 (113d) 自动满足. 因为若 $(A_0,A,\phi)$ 满足 (113a)--(113c), 则 $u=\partial^j A_j$ 满足 $\Box u=|\phi|^2u$, 而若 (114) 与 (115) 成立, 则 $u|_{t=0}=\partial_t u|_{t=0}=0$.
于是我们只需方程 (113a)--(113c). 其中第一个为椭圆方程, 较易处理, 故将其排除在我们的模型方程之外. 剩余两个方程右端含有三类项:
\begin{itemize}
\item 涉及 $A_0$ 的"椭圆项", 统称 $\mathcal{E}$;
\item $A_j$ 与 $\phi$ 中的三次项, 统称 $\mathcal{C}$;
\item 具零形式结构的二次项.\end{itemize}

归入后一类的项为 $-\Im(\phi\,\overline{\partial_j\phi})$ 与 $-2\mathrm{i}A^j\partial_j\phi$. 我们现在揭示这些表达式内蕴的零形式结构(源于 Coulomb 条件).
将 $\phi$ 分解为实部与虚部: $\phi=u+\mathrm{i}v$. 则
\begin{equation}
-\Im(\phi\,\overline{\partial_j\phi})=u\partial_j v-v\partial_j u,\end{equation}
于是 (113b) 作为(依赖于时间的)一次形式方程在 $\mathbb{R}^n$ 上可写为
\begin{equation}
\Box A=u\,dv-v\,du+\mathcal{C}+\mathcal{E}.\end{equation}
两边同时作用外微分 $d$:
\begin{equation}
\Box(dA)=2\,du\wedge dv+d\mathcal{C}.\end{equation}
而
\begin{equation}
du\wedge dv=\frac12 Q_{jk}(u,v)\,dx^j\wedge dx^k,\end{equation}
由此
\begin{equation}
\Box F_{jk}=Q_{jk}(u,v)+\partial\mathcal{C}.\end{equation}
Coulomb 规范条件推出 $\partial^k F_{jk}=-\Delta A_j$, 于是
\begin{equation}
-\Delta\Box A_j=\partial^k Q_{jk}(u,v)+\partial^2\mathcal{C}.\end{equation}
从而在 Riesz 算子意义下,
\begin{equation}
\Box A=D^{-1}Q(\Re\phi,\Im\phi)+\mathcal{C},\end{equation}
其中 $Q$ 为零形式 $Q_{jk}$ 的某个线性组合. 由于三次项 $\mathcal{C}$ 更易估计, 我们将其排除在模型问题之外.
现考虑方程 (113c). 分离实部与虚部, 有
\begin{equation}
\begin{aligned}
\Box u&=2A\cdot\nabla v+\mathcal{C}+\mathcal{E},\\
\Box v&=-2A\cdot\nabla u+\mathcal{C}+\mathcal{E}.\end{aligned}
\end{equation}
(此处我们将 $A$ 借指标提升视为向量场; $\nabla$ 表示空间变量的梯度. )我们断言, 项 $A\cdot\nabla u$ 与 $A\cdot\nabla v$ 具有零形式结构, 这是因为 $A$ 无散度(由 Coulomb 条件). 设 $B_{jk}$ 为方程
\begin{equation}
\Delta B_{jk}=\partial_j A_k-\partial_k A_j
\end{equation}
的唯一解(在适当的正则性假设下). 由 Coulomb 条件,
\begin{equation}
\Delta\partial^j B_{jk}=\Delta A_k,\qquad\text{即}\quad \partial^j B_{jk}=A_k.\end{equation}
于是
\begin{equation}
A\cdot\nabla u=\partial^j B_{jk}\partial^k u=\frac12 Q_{jk}(u,B^{jk}).\end{equation}
上述关于 $u=\Re\phi$ 与 $v=\Im\phi$ 的方程因而可重写为
\begin{equation}
\begin{aligned}
\Box\Re\phi&=Q_{jk}(\Im\phi,B^{jk})+\mathcal{C}+\mathcal{E},\\
\Box\Im\phi&=Q_{jk}(B^{jk},\Re\phi)+\mathcal{C}+\mathcal{E}.\end{aligned}
\end{equation}
但由 (117), $B$ 在 Riesz 算子意义下形如 $D^{-1}A$. 将此与 (116) 结合, 并处处舍去项 $\mathcal{C}$ 与 $\mathcal{E}$, 我们得到形如 ("MKG") 的一个系统, 此即我们的 (MKG) 模型.
\section*{5.0.3 Yang--Mills 方程}

设 $G$ 为经典紧 Lie 矩阵群之一(如 $SO(k,\mathbb{R})$ 或 $SU(k,\mathbf{C})$), $\mathfrak{g}$ 为其 Lie 代数. 未知量为 $\mathbb{R}^{1+n}$ 上取值于 $\mathfrak{g}$ 的一次形式 $A_\mu dx^\mu$. 相应的协变导数为
\begin{equation}
D_\mu H=\partial_\mu H+[A_\mu,H],\end{equation}
其中 $H$ 为任意取值于 $\mathfrak{g}$ 的张量场, $[\cdot,\cdot]$ 为矩阵交换子. 曲率(曲率为取值于 $\mathfrak{g}$ 的二次形式)为
\begin{equation}
F_{\mu\nu}=\partial_\mu A_\nu-\partial_\nu A_\mu+[A_\mu,A_\nu].\end{equation}
作用量为
\begin{equation}
\mathcal{L}[A_\mu]=-\frac14\int\langle F_{\mu\nu},F^{\mu\nu}\rangle\,dt\,dx,\end{equation}
其中 $\langle\cdot,\cdot\rangle$ 为 $\mathfrak{g}$ 上继承自环境空间的内积(例如 $SO(k,\mathbb{R})$ 嵌入 $\mathbb{R}^{k^2}$, 其 Lie 代数可视为后者的子空间). Euler--Lagrange 方程为
\begin{equation}
D^\nu F_{\mu\nu}=0.\end{equation}
设 $O$ 为 $\mathbb{R}^{1+n}$ 上取值于 $G$ 的函数. 考虑规范变换 $A_\mu\to\widetilde A_\mu$, 由
\begin{equation}
\widetilde A_\mu=O A_\mu O^{-1}-\partial_\mu O\,O^{-1}
\end{equation}
给出. 计算表明曲率随之变换为
\begin{equation}
\widetilde F_{\mu\nu}=O F_{\mu\nu}O^{-1}.\end{equation}
记 $\widetilde D_\mu$ 为对应于 $\widetilde A_\mu$ 的协变导数, 则
\begin{equation}
\widetilde D^\nu\widetilde F_{\mu\nu}=O D^\nu F_{\mu\nu}O^{-1},\end{equation}
故 (119) 在规范变换下不变. 于是我们拥有规范自由度, 并可对 $A_\mu$ 施加规范条件.
\paragraph*{(YM) 的 Coulomb 规范.} 相对 Coulomb 条件 $\partial^j A_j=0$, (119) 取形式(见 \cite{klai34})
\begin{gather}
\Delta A_0=2[\partial^j A_0,A_j]+[A^j,\partial_t A_j]+[A^j,[A_0,A_j]],\\
\Box A_j+\partial_t\partial_j A_0=-2[A^k,\partial_k A_j]+[A^k,\partial_j A_k]+[\partial_t A_0,A_j]+2[A_0,\partial_t A_j]\nonumber\\
\qquad-\,[A_0,\partial_j A_0]-[A^k,[A_k,A_j]]+[A_0,[A_0,A_j]],\\
\partial^j A_j=0.\end{gather}
遗憾的是, 假设存在整体 Coulomb 规范, 会对初数据施加一个限制性的小性假设. 在 \cite{klai34} 中, 这一困难借助局部论证得以解决. 依 \cite{klai47}, 我们忽略这一复杂之处, 而从系统 (119) 导出我们的模型方程.
与 (MKG) 的讨论一样, (119c) 化为初数据上的一个约束. 关于 $A_0$ 的方程是椭圆的, 故我们忽略它. 至于 (119b), 我们只保留右端的前两项, 因为其余项要么涉及 $A_0$(对其我们期望得到比 $A_j$ 更好的估计), 要么为三次项.
现将 (119b) 写为(依赖于时间的、取值于 $\mathfrak{g}$ 的)$\mathbb{R}^n$ 上一次形式的方程(忽略右端除前两项外的一切):
\begin{equation}
\Box A+d(\partial_t A_0)=S+T,\end{equation}
其中 $A=A_j dx^j$, $S=-2[A^k,\partial_k A_j]dx^j$, $T=[A^k,\partial_j A_k]dx^j$. 两边同时作用外微分 $d$:
\begin{equation}
\Box dA=dS+dT.\end{equation}
设 $B$ 为 (117) 所确定的二次形式(此处取值于 $\mathfrak{g}$). 于是
\begin{equation}
\Box\Delta B_{jk}=\partial_j S_k-\partial_k S_j+\partial_j T_k-\partial_k T_j.\end{equation}
由 (118) 推得
\begin{equation}
-\Box\Delta A_j=\partial^k(\partial_j S_k-\partial_k S_j+\partial_j T_k-\partial_k T_j),\end{equation}
故就频率空间中的估计而言, 我们可用下式替代 (119b):
\begin{equation}
\Box A=S+D^{-1}dT.\end{equation}
剩下要辨明隐藏于 $S$ 与 $dT$ 中的零形式结构. 首先,
\begin{align*}
S_j&=[\partial_k B^{kl},\partial_l A_j]=\frac12[\partial_k B^{kl}-\partial_k B^{lk},\partial_l A_j]\\
&=\frac12\bigl([\partial_k B^{kl},\partial_l A_j]-[\partial_l B^{kl},\partial_k A_j]\bigr)\\
&=\frac12\bigl(\partial_k B^{kl}\partial_l A_j-\partial_l B^{kl}\partial_k A_j+\partial_k A_j\partial_l B^{kl}-\partial_l A_j\partial_k B^{kl}\bigr),\end{align*}
故矩阵 $S_j$ 的每一项都是形如 $Q_{kl}(A,B)$ 的项的线性组合, 其中 $A$ 与 $B$ 表示 $A_j$ 与 $B_{kl}$ 的任意两个分量. 但由 (118), 我们可用 $D^{-1}A$ 替代 $B$. 概略地,
\begin{equation}
S\sim Q(A,D^{-1}A).\end{equation}
现考虑一次形式 $T$. 我们计算:
\begin{align*}
(dT)_{jk}&=\partial_j[A^l,\partial_k A_l]-\partial_k[A^l,\partial_j A_l]\\
&=[\partial_j A^l,\partial_k A_l]-[\partial_k A^l,\partial_j A_l].\end{align*}
于是矩阵 $(dT)_{jk}$ 的每一项都是形如 $Q_{jk}(A,A')$ 的项的线性组合, 其中 $A$ 与 $A'$ 表示 $A_l$($1\leq l\leq n$)的任意两个分量. 将此与 (121) 及 (120) 结合, 我们即得 Yang--Mills 方程的模型 ("YM").
